% Generated draft: requires mathematical and source review.
\documentclass[11pt]{article}
\usepackage[margin=25mm]{geometry}
\usepackage{fontspec}
\setmainfont{Latin Modern Roman}
\usepackage{amsmath,amssymb,mathtools,unicode-math}
\setmathfont{Latin Modern Math}
\usepackage{enumitem,needspace,xurl,hyperref}
\setlength{\parindent}{0pt}
\setlength{\parskip}{5pt}
\setlength{\emergencystretch}{3em}
\newcommand{\N}{\mathbb{N}}
\newcommand{\Z}{\mathbb{Z}}
\newcommand{\Q}{\mathbb{Q}}
\newcommand{\R}{\mathbb{R}}
\newcommand{\C}{\mathbb{C}}
\newcommand{\F}{\mathbb{F}}
\newcommand{\A}{\mathbb{A}}
\newcommand{\PP}{\mathbb{P}}
\newcommand{\E}{\mathbb{E}}
\newcommand{\eps}{\varepsilon}
\newcommand{\dd}{\,\mathrm{d}}
\begin{document}
\section*{UnsolvedMath ID: 3100017}
\textbf{Source problem code:} catalogue entry
\subsection{Definitions and mathematical statement}
Let $G=(V,E)$ be a finite simple undirected graph. For disjoint $S,T\subseteq V$, let $e_G(S,T)$ count edges with one endpoint in each set, and define
\[
 \operatorname{disc}(G)=\max_{S\cap T=\varnothing}
 \left|\frac{|S||T|}{2}-e_G(S,T)\right|.
\]
The reference density is exactly $1/2$, not the actual edge density of $G$.

An induced copy of a finite graph $H$ is an injective map from $V(H)$ into $V(G)$ preserving both adjacency and nonadjacency. Chung and Graham ask how large $\operatorname{disc}(G)$ must be if $G$ contains no induced copy of $H$.

For each pair $(n,H)$ for which the class is nonempty, the sharp extremal quantity can be written
\[
 d_H(n)=\min\{\operatorname{disc}(G):|V(G)|=n,
                          \ G\text{ has no induced copy of }H\}.
\]
The problem is to determine useful or optimal lower bounds on $d_H(n)$, with the dependence on $H$ made explicit. This formulation allows $H$ to be fixed or to vary with $n$; the source does not state which regime, whether an asymptotic order or an optimal constant is wanted, or a proposed numerical bound. For exceptional pairs with no admissible graph, no finite minimum is silently assigned. Ordinary subgraph exclusion is not substituted for induced exclusion.
\subsection{Short English statement}
Quantify how far an induced-$H$-free graph must be from having half the possible edges across every pair of disjoint vertex sets. The source leaves the intended parameter regime and sharpness target unspecified. Its assessment mentions quantitative discrepancy work of Erdős--Spencer and Chung--Graham and links to Erdős--Hajnal/pseudorandomness, but supplies no precise theorem or verified full answer. The October 2020 open label and later unsuccessful resolution check are dated source evidence.
\subsection{Sources}
\begin{itemize}
\item UnsolvedMath permanent ID 3100017, source code catalogue entry, Chung--Graham's induced-subgraph discrepancy question. Registry curation and metadata: Ulam AI and UnsolvedMath Contributors, CC BY 4.0; updated 2026-08-06.
\item Cooper, \emph{Combinatorial Problems I Like}, Graphs, Hypergraphs, Set Systems, Designs, source-order bullet 17, October 2020 snapshot: \url{https://people.math.sc.edu/cooper/combprob.html}, supplied HTML extraction accessed 2026-07-29. Later resolution was not checked item by item; redistribution permission was not established in the supplied metadata.
\item Erdős--Spencer, Chung--Graham, and the Erdős--Hajnal connection are named by the supplied partial-progress/OPEN-TRIAGE assessment without complete references or theorem statements. No missing citation is inferred.
\end{itemize}
\end{document}
